\section{Related Work}
\label{sec:related-work}

\paragraph{Adaptive geometry and variational reduction.}
Natural gradient, K-FAC, AdaGrad, Adam, Shampoo, SOAP, and related methods use
Fisher, moment, Kronecker, or tensor states as changing positive update
operators \cite{amari1998natural,martens2015kfac,duchi2011adagrad,
kingma2014adam,gupta2018shampoo,morwani2024shampoo,vyas2025soap}.
Riemannian optimization provides the coordinate-free gradient language
\cite{absil2008optimization,boumal2023introduction}.  Partial minimization,
marginal functions, envelope sensitivity, and variable projection are
classical \cite{rockafellar1970convex,rockafellar1998variational,
bonnans2000perturbation,golub1973variableprojection}.  In fixed affine
parameter coordinates, the Schur complement is the reduced coordinate Hessian
of a marginal objective.  Global
first-order complexity for geodesically convex objectives on Hadamard
manifolds is established in \cite{zhang2016firstorder}, and metric projection
and convex optimization in Hadamard spaces are developed systematically by
\cite{bacak2014convex}.  Our use of these tools is tied to the action
constraint: the global fiber coordinates expose an exact logarithmic
residual, a current-sublevel curvature majorant attained by the action family,
an active realization, and posterior controller certificates.

\paragraph{AIRM approximation and totally geodesic submanifolds.}
The affine-invariant SPD geometry is classical
\cite{bhatia2007positive,pennec2006riemannian}.  AIRM approximation from
special sets, best approximation in geodesic SPD submanifolds, and general
total-geodesy criteria and projections are studied by
\cite{bhatia2014approximation}, \cite{lim2004best}, and
\cite{tumpach2024totally}; conic geodesic optimization supplies a broader
algorithmic SPD setting \cite{sra2015conic}.  Algebraic identities and
definiteness conditions for Hermitian solutions of \(AX=B\) are studied by
\cite{tian2013hermitian}.  These results supply feasibility and block-Schur
algebra, while Hadamard projection theory supplies uniqueness once a fiber is
known to be closed and geodesically convex.  \cite{lim2004best} additionally
gives explicit formulas for selected involutive block models, including
low-dimensional block cases.

The comparison is result-specific.  Tian's matrix-equation analysis supplies
algebraic existence, definiteness, and solution identities for Hermitian
equations; our \(A^\top B\succ0\) condition is used at that feasibility level,
whereas the determinant-one AIRM minimizer and its iteration are variational
objects.  Lim studies best approximation in Riemannian geodesic SPD
submanifolds and gives explicit formulas for selected involutive block models;
we use the same nearest-point geometry only after proving that every
action-indexed fiber is closed and totally geodesic.  Tumpach and Larotonda
give general total-geodesy criteria and structural results for SPD
submanifolds; our contribution identifies the full \(PA=B\) family, constructs
its global analytic bundle over \(B\), and derives the normal equation.  None
of these three results supplies the action residual iteration, its
current-sublevel rate and posterior certificates, the active-compatible
\(r\le2m\) realization, or the rank-\(d-1\) normalized recovery law.  Conversely,
projection uniqueness and the general variational/Schur principles are used as
established ingredients rather than claimed as new.  The theorem-level map in
\cref{app:closest-comparison} records these boundaries.

\paragraph{Secant updates and positive-definite completion.}
Inverse-Hessian multi-secant equations have the form \(PY=S\)
\cite{schnabel1983multiple}; compatibility with positive definiteness and
modern block or grouped updates have been studied separately
\cite{passy1984secant,gao2018block,boutet2022grouped}.  Classical least-change
and variational BFGS/DFP updates use weighted norms or log-determinant
divergences \cite{dennis1979leastchange,fletcher1991variational}; Bregman and
weighted secant extensions use related selection principles
\cite{kanamori2013bregman,gratton2015weighted}.  Bounded and sparse
quasi-Newton updates and positive-definite completion are treated in
\cite{calamai1987bounds,yamashita2007sparse,dai2014sparse,grone1984positive}.

Our action formulation gives a global analytic trivialization of
\(P\mapsto PA\), identifies its determinant-one fibers, and proves their total
geodesy before applying the nearest-point principle.  It then yields a
globally convergent residual solver and the sharp completion rank for
shape-normalized nested actions.  The rank-\(d-1\) theorem uses a supplied
determinant scale, making its information model different from ordinary
scale-free secant observations.  Condition-constrained SPD approximation and
preconditioning have separate literatures
\cite{tanaka2014condition,marechal2009optimizing,lu2011minimizing,
gao2023scalable,qu2024optimal,dogan2025geodesicpreconditioning}.  A
theorem-level comparison is given in \cref{app:closest-comparison}.

\paragraph{Scope.}
The companion information-compression map \(P\mapsto A^\top PA\) has different
fibers \informationgeometrycite.  The present paper is self-contained and
treats the full-SPD exact-action problem.  Structured controller families,
noisy secants, unknown scalar gauges, historical state estimation, and regret
lead to distinct constrained-optimization and statistical problems.
